%% ── 12. APRENDIZAJES Y LECCIONES ─────────────────────────────────────────────
\chapter{Aprendizajes y Lecciones}

Este capítulo documenta los aprendizajes más significativos extraídos del
ciclo completo de desarrollo y evaluación de ContractIA: desde los primeros
experimentos en notebook hasta el MVP desplegado en producción. Se organizan
en cuatro categorías: técnicos (arquitectura e implementación), de
infraestructura y operaciones, metodológicos (proceso de desarrollo), y de
negocio (adopción e impacto institucional).

\section*{Aprendizajes Técnicos}

\begin{enumerate}

    \item \textbf{Los agentes especializados superan ampliamente al agente único.}

    La descomposición del análisis en tres agentes con roles diferenciados
    (Jurista, Auditor, Cronista) produjo $+77\%$ más hallazgos respecto al
    agente único de la v1.0 (119 vs.\ 35 en el contrato piloto), sin
    degradar la precisión. El beneficio no proviene solo de la paralelización,
    sino de que cada agente recibe un prompt ajustado a su dominio semántico:
    el Jurista responde mejor a preguntas de cumplimiento normativo que un
    agente generalista que también debe detectar plazos y contradicciones.

    \textit{Principio derivado:} en tareas de análisis documental complejo,
    la especialización de agentes por tipo de tarea genera más valor que
    aumentar la complejidad del prompt de un único agente.

    \item \textbf{El RAG híbrido es esencial para análisis legal de precisión.}

    La búsqueda semántica pura (FAISS) falla en la recuperación de referencias
    legales específicas como \texttt{Artículo 18.3 del D.Leg.\ 1362} o
    \texttt{cláusula 6.2.b} del propio contrato, ya que estos identificadores
    no tienen un vecino semántico cercano en el espacio de embeddings.
    El componente léxico BM25 resuelve este problema directamente.
    El Reranker Cohere actúa como un tercer filtro que prioriza la relevancia
    contextual sobre la puntuación bruta de recuperación.

    \textit{Principio derivado:} para dominios con vocabulario técnico-legal
    preciso, el RAG híbrido (léxico + semántico + reranker) no es un lujo
    sino un requisito de precisión.

    \item \textbf{GraphRAG actúa como amplificador o filtro según el modelo LLM.}

    Con Gemini~2.5~Pro, el grafo amplifica los hallazgos al identificar
    relaciones \texttt{contradice} entre cláusulas distantes. Con
    Gemini~3.1~Pro~Preview, el grafo más denso (477 nodos vs.\ 208) reduce
    los hallazgos porque el modelo usa la riqueza relacional para descartar
    tensiones textuales que no alcanzan el umbral de inconsistencia real.
    Ninguna de las dos conductas es incorrecta: reflejan el trade-off
    sensibilidad/precisión del modelo subyacente.

    \textit{Principio derivado:} el efecto de GraphRAG no puede evaluarse
    de forma aislada; siempre está mediado por el comportamiento del LLM
    que construye y consume el grafo.

    \item \textbf{\texttt{ContextVar} no propaga en \texttt{run\_in\_executor}.}

    Python propagates \texttt{contextvars.ContextVar} en coroutines
    asíncronas, pero no en threads del \texttt{executor} de asyncio.
    Esto causó que el \texttt{audit\_id} del trabajo en curso no estuviera
    disponible en los módulos de segmentación y GraphRAG que corren en
    \texttt{run\_in\_executor}, rompiendo el sistema de logging en producción.
    La solución fue eliminar el \texttt{ContextVar} y pasar \texttt{audit\_id}
    explícitamente como parámetro en todas las funciones del pipeline.

    \textit{Principio derivado:} en arquitecturas que mezclan \texttt{asyncio}
    y \texttt{ThreadPoolExecutor}, evitar el estado implícito por
    \texttt{ContextVar}; preferir paso explícito de parámetros de contexto.

    \item \textbf{Los esquemas Pydantic son el mecanismo de contrato entre el LLM y el sistema.}

    Sin validación Pydantic, la salida JSON del LLM es impredecible: campos
    faltantes, tipos incorrectos, texto libre donde se espera un enum.
    Los esquemas estrictos (\texttt{model\_config = ConfigDict(strict=True)})
    forzaron al LLM a respetar el contrato de datos en $> 98\%$ de las
    llamadas, y permitieron detectar y reintentar el $< 2\%$ restante
    sin propagar datos corruptos al consolidador.

    \textit{Principio derivado:} en sistemas LLM en producción, la validación
    estructurada de la salida no es opcional; es la interfaz entre el
    componente no determinista y el sistema determinista.

    \item \textbf{El módulo anti-TOC es tan crítico como el RAG.}

    En las primeras iteraciones, sin anti-TOC, el índice de contenido del
    contrato (TOC de 12 páginas con referencias del tipo ``6.2.b ..... 87'')
    era indexado como fragmentos de cláusulas. El RAG recuperaba estos
    fragmentos y los agentes generaban hallazgos falsos sobre ``obligaciones''
    que eran en realidad entradas del índice. El módulo anti-TOC basado en
    densidad de referencias numéricas y longitud mínima de fragmento
    eliminó este problema completamente.

    \textit{Principio derivado:} en contratos legales, la calidad de la
    segmentación upstream determina el techo de precisión de todo el
    pipeline downstream.

\end{enumerate}

\section*{Aprendizajes de Infraestructura y Operaciones}

\begin{enumerate}

    \item \textbf{Cloud Run no soporta SSE nativo sin proxy.}

    Las respuestas de tipo Server-Sent Events (SSE) son bufferizadas por
    el proxy de Cloud Run antes de llegar al cliente, lo que hace que el
    cliente no reciba actualizaciones en tiempo real sino un único mensaje
    al final. Sin desplegar un proxy nginx por delante, SSE es inviable
    en Cloud Run. La solución adoptada fue reemplazar SSE con polling
    cada 5--10 segundos desde el cliente, que es más simple, sin estado
    en el servidor, y compatible con cualquier proxy intermedio.

    \textit{Principio derivado:} en entornos serverless (Cloud Run, Lambda,
    Cloud Functions), los patrones de streaming requieren verificación
    previa con la capa de red del proveedor.

    \item \textbf{Workload Identity Federation elimina la necesidad de claves de servicio.}

    Usar un \textit{Service Account Key} JSON almacenado en GitHub Secrets
    es un riesgo de seguridad: la clave puede filtrarse, no tiene rotación
    automática, y es difícil de revocar selectivamente. Migrar a Workload
    Identity Federation (WIF) permitió que GitHub Actions autentique contra
    GCP sin almacenar ninguna clave, reduciendo la superficie de ataque del
    CI/CD a cero secretos de largo plazo.

    \textit{Principio derivado:} en pipelines CI/CD que despliegan a la nube,
    WIF (o equivalente federado) es siempre preferible a claves de servicio
    almacenadas.

    \item \textbf{Los análisis de larga duración requieren diseño asíncrono desde el inicio.}

    Diseñar el endpoint \texttt{/audit} como síncrono hubiera limitado
    el sistema a contratos que completan en $< 60$ segundos (timeout de
    Cloud Run). Diseñarlo como asíncrono desde v9.5 (dispara
    \texttt{BackgroundTask}, devuelve \texttt{job\_id}) permitió escalar
    a análisis de 75+ minutos sin cambiar la interfaz del cliente. El
    sistema de polling y notificación por email fue una consecuencia natural
    de esta decisión arquitectónica, no un parche posterior.

    \textit{Principio derivado:} para tareas de larga duración, el patrón
    \textit{fire-and-poll} debe ser la primera opción de diseño, no la última.

    \item \textbf{Secret Manager justifica su costo operativo desde el primer secreto.}

    Gestionar credenciales en variables de entorno del Dockerfile o en
    archivos \texttt{.env} es frágil, difícil de auditar y propenso a
    fugas en logs. Centralizar todos los secretos en Cloud Secret Manager
    con acceso por service account eliminó los secretos del código base,
    del historial de Git y de los logs de Cloud Run, y habilitó la
    auditoría de acceso por secreto en Cloud Audit Logs.

\end{enumerate}

\section*{Aprendizajes Metodológicos}

\begin{enumerate}

    \item \textbf{Desplegar en producción desde etapas tempranas revela problemas invisibles en notebook.}

    El prototipo en notebook (v1.0) procesó el Contrato A en 47 minutos
    con 35 hallazgos. El MVP en producción (v9.8) procesa el Contrato A
    en 75 minutos con 119 hallazgos. La diferencia no es solo de mejoras
    arquitectónicas: el entorno de Cloud Run expuso problemas de
    \textit{throttling} de la API de Vertex AI (429 en horas pico), latencias
    de red variables hacia Cohere, y tiempos de escritura a Cloud SQL que
    el notebook local nunca mostró. Sin producción temprana, estos problemas
    habrían emergido solo en la fase de lanzamiento.

    \item \textbf{Un único contrato de referencia fijo acelera la evaluación iterativa.}

    Usar el mismo contrato (Contrato A) como benchmark constante
    en todas las iteraciones permitió comparaciones directas y objetivas:
    ``v8.8 detectó 89 hallazgos en 65 min; v9.0 detectó 103 en 72 min''.
    Sin un benchmark fijo, las mejoras iterativas habrían sido anecdóticas
    y difíciles de comunicar a stakeholders.

    \item \textbf{La evaluación comparativa de modelos genera conocimiento metodológico no anticipado.}

    Correr Gemini~2.5~Pro y Gemini~3.1~Preview sobre el mismo contrato
    no era un requisito del MVP; fue una decisión experimental. El hallazgo
    de la ``paradoja grafo-hallazgos'' (mayor grafo $\neq$ más hallazgos)
    no solo es valioso para ContractIA sino para el diseño de cualquier
    sistema que use GraphRAG con LLMs de diferente capacidad de razonamiento.
    Los experimentos comparativos, aunque costosos, son una fuente de
    conocimiento sistematizable.

    \item \textbf{El Lean Startup reduce el riesgo en proyectos de IA aplicada.}

    Comenzar con la hipótesis más simple (¿puede un LLM detectar una
    inconsistencia en un contrato?) y validarla antes de construir el
    pipeline completo redujo el riesgo de invertir semanas en una
    arquitectura que luego resultara inviable. Cada iteración resolvió
    un único problema, medido contra el benchmark fijo, antes de agregar
    nueva complejidad.

\end{enumerate}

\section*{Aprendizajes de Negocio e Institucional}

\begin{enumerate}

    \item \textbf{La reducción de 320 horas a 1 hora es el argumento de adopción más poderoso.}

    En las conversaciones con especialistas de ProInversión, ningún argumento
    técnico (precisión, cobertura, trazabilidad) tuvo el impacto comunicativo
    de este número. ``320 horas de trabajo de 5 especialistas reducidas a
    75 minutos'' traduce el valor técnico al idioma de la gestión institucional.
    Para la adopción de herramientas de IA en el sector público, el ROI
    expresado en tiempo-hombre es más persuasivo que cualquier métrica técnica.

    \item \textbf{El reporte ejecutivo es tan importante como los hallazgos.}

    El sistema podría detectar 200 inconsistencias y aun así no ser adoptado
    si los resultados no pueden ser consumidos eficientemente. Los directivos
    necesitan un resumen ejecutivo de 2--3 páginas con los 10 hallazgos
    de mayor riesgo y su impacto potencial. Los especialistas necesitan
    el reporte técnico completo con citas textuales y referencias normativas.
    Invertir en los dos formatos de salida es tan importante como invertir
    en la precisión del análisis.

    \item \textbf{La trazabilidad es el requisito no negociable en entornos institucionales.}

    Sin ubicación exacta (capítulo, sección, número de cláusula), fragmento
    de evidencia literal y referencia normativa, un hallazgo generado por
    IA no puede ser usado en un proceso formal de auditoría contractual.
    Los especialistas necesitan poder ir directamente al fragmento del
    contrato que origina el hallazgo para evaluarlo; sin esa capacidad,
    el hallazgo no tiene valor operativo. Esta lección moldeó el diseño
    del esquema Pydantic de cada agente desde la v5.

    \item \textbf{La IA como apoyo, no como reemplazo, reduce la resistencia institucional.}

    Presentar ContractIA como un sistema de ``screening previo'' que
    reduce el tiempo de preparación del especialista, en lugar de
    ``sistema que reemplaza la revisión legal'', eliminó las objeciones
    más comunes de los usuarios potenciales. El framing correcto del
    producto es tan importante como el producto mismo para la adopción
    en entidades con alta sensibilidad hacia la responsabilidad profesional.

    \item \textbf{Los hallazgos más valiosos son los que el especialista no habría encontrado.}

    Los 5 hallazgos confirmados por ambos modelos (sección 10) son
    significativos, pero los más impactantes durante las demos fueron
    los hallazgos de referencias cruzadas entre cláusulas no adyacentes
    (por ejemplo, la contradicción entre 4.7 y 4.13--4.14): inconsistencias
    que un revisor leyendo secuencialmente rara vez detecta, pero que el
    sistema encuentra sistematicamente al analizar el grafo de relaciones.
    Este tipo de hallazgo ``imposible de encontrar manualmente'' es el
    argumento más fuerte para la adopción del sistema.

\end{enumerate}

\section*{Síntesis: Principios para Sistemas LLM en Producción}

Los aprendizajes anteriores se condensan en los siguientes principios
aplicables a cualquier sistema de análisis documental con LLMs:

\begin{table}[H]
\centering
\caption{Principios derivados del desarrollo de ContractIA}
\label{tab:principios}
\begin{tabular}{p{1cm} p{4.5cm} p{8cm}}
\toprule
\textbf{\#} & \textbf{Principio} & \textbf{Aplicación} \\
\midrule
1 & Especialización $>$ generalización &
    Múltiples agentes con prompts diferenciados generan más valor que un
    agente generalista con prompt largo. \\
2 & RAG híbrido para dominios técnicos &
    BM25 + semántico + reranker para dominios con vocabulario preciso. \\
3 & Validación estructurada de salida &
    Pydantic o equivalente; nunca confiar en JSON libre del LLM. \\
4 & Calidad de segmentación = techo de precisión &
    El pipeline es tan preciso como el módulo de extracción upstream. \\
5 & Asincronía desde el inicio &
    Tareas $> 60$ s deben ser \textit{fire-and-poll} desde el diseño. \\
6 & Producción temprana &
    Desplegar en Cloud Run desde v5 expone problemas invisibles en notebook. \\
7 & Benchmark fijo &
    Un único caso de referencia fijo permite comparaciones iterativas objetivas. \\
8 & ROI en tiempo-hombre &
    Para adopción institucional, el argumento es el tiempo reducido, no la
    métrica técnica. \\
\bottomrule
\end{tabular}
\end{table}